\chapter{ХЭЛЭЛЦҮҮЛЭГ БА ЦААШДЫН АЖИЛ}
\section{Үр дүнгийн хэлэлцүүлэг}
\textbf{Илрүүлэгч.} Дугаарын геометр улс орноос үл хамаардаг тул 8\,823 олон улсын зураг дээр сургасан YOLO11n монгол дугаарыг шууд өндөр нарийвчлалтай олж (mAP$_{50}$ 0.98), 48 монгол зураг нэмэхэд хайрцгийн нарийвчлал (mAP$_{50\text{--}95}$ 0.69$\to$0.83) л сайжирсан. Энэ нь илрүүлэгчид монгол зураг харьцангуй цөөн байхад ч хангалттай, харин уншигчид бодит монгол өгөгдөл шийдвэрлэх ач холбогдолтой гэсэн дүгнэлтэд хүргэнэ.

\textbf{Бодит өгөгдлийн ач холбогдол.} Зөвхөн синтетик 42 хувь, 74 бодит мөр нэмэхэд 84 хувь, 7\,392 бодит тайралт нэмэхэд 97 хувь болсон нь (Хүснэгт~\ref{tab:ablation}) синтетик нь тэмдэгтийн хэлбэрийг зааж, бодит өгөгдөл домэйны онцлогийг заадаг гэсэн~\cite{jaderberg2014,ruder2022}-ын дүгнэлтийг давтаж байна. Монгол дугаарт текстээр шошголсон нээлттэй өгөгдөл байхгүй гэж үзэж байсан ч GitHub репо дотроос MLUB тэмцээний өгөгдөл олдсон нь ажлын чиглэлийг өөрчилсөн.

\textbf{Яагаад төгсгөл хүртэлх үр дүн доогуур вэ.} 1\,491 тайралт дээрх 96.8 хувь ба 19 бүтэн дүрстэй зураг дээрх 73.7 хувийн зөрүү таван шалтгаантай: (1) текстээр шошголсон бодит монгол өгөгдөл 7\,392 зураг, 5\,575 ялгаатай дугаар байгаа нь хятад (CCPD 300 мянга) зэрэгтэй харьцуулахад бага; (2) үсгийн тархалт тэнцвэргүй --- аймгийн кодоос болж У, Б, Н, А мянгаар байхад Л, К, М, И, Т, Г серийн үсэг 100--300 удаа л байгаа ба алдсан дөрвөн дугаар яг эдгээр үсэг дээр алдсан; (3) домэйн зөрүү --- MLUB замын камерын тайралт бол GitHub гар утасны бүтэн машины зураг, бохир, цастай, налуу, энэ домэйноос сургалтад 74 мөр л орсон; (4) илрүүлэгчийн хайрцаг үнэн хайрцгаас зөрөхөд мөр хуваалт өөрчлөгдөж 89.5-аас 77.8 хувь болдог; (5) 19 дугаар дээр нэг алдаа 5.3 нэгж хувь тул статистикийн хувьд хангалтгүй (95 хувийн Wilson интервал 51--88).

\textbf{Олон кадрын санал.} Симуляцид саналын ашиг бага (73.7$\to$78.9, MLUB-д өөрчлөлтгүй) байсан нь нэг зургийн хувилбарууд дээр алдаа системтэй учраас; бодит бичлэгт нэг кадрын санал нийцэл 0.3--0.8 байхад санал нэгдсэн дугаар тогтвортой гарч байгаа нь (Зураг~\ref{fig:tracks}) бодит нөхцөлд саналын ач холбогдол илүү гэдгийг харуулна. Margin ба нийцэл хоёр нь найдваргүй мөрийг ялгах тул хаалга нээх шийдвэрт заавал ашиглах хэрэгтэй.

\section{Хязгаарлалт}
\begin{itemize}[leftmargin=2em]
\item Бодит зураг бүгд өдрийн, өвлийн, гар утас эсвэл замын камерын; шөнө, бороо, гэрлийн гялбаа, хаалганы доод өнцгийг зөвхөн аугментаци дуурайсан.
\item Дипломат (улаан, ДК), цэргийн (ЦАБ), хилийн (ХЦА), чиргүүл, мотоциклийн дугаар байхгүй.
\item Хоёр мөрт хуваалт тогтмол 55/55 хувь; хүчтэй налуу (roll) үед мөр холилдоно.
\item MLUB хэлбэрийн шошго автомат (3 хувь орчим алдаатай); тестийн тоо үүгээр бага зэрэг дутуу үнэлэгдэнэ.
\item Төгсгөл хүртэлх тест 19 дугаартай; Roboflow-ийн 984 зураг API түлхүүргүй тул ашиглагдаагүй.
\item Бичлэгийн үр дүн үнэн шошгогүй тул чанарын.
\end{itemize}

\section{Явцын төлөв ба үлдсэн ажил}
\begin{table}[H]\centering\caption{Гүйцэтгэсэн ба үлдсэн ажил}\label{tab:todo}\small
\begin{tabular}{p{9cm}l}\toprule
Ажил & Төлөв\\\midrule
Судалгааны тойм, монгол дугаарын бүтэц, эх үүсвэрийн судалгаа & гүйцэтгэсэн\\
GitHub, HuggingFace, MLUB өгөгдөл бэлтгэх, текст сэргээх, давхцалгүй хуваалт & гүйцэтгэсэн\\
Синтетик дугаар үүсгэгч (90\,021 мөр) & гүйцэтгэсэн\\
YOLO11n хоёр үе шаттай сургалт, mAP$_{50}$ 0.992 & гүйцэтгэсэн\\
CRNN+CTC дөрвөн үе шаттай сургалт, 96.8 хувь & гүйцэтгэсэн\\
Дугаарын дүрэм, аймгийн код, олон кадрын жинлэсэн санал & гүйцэтгэсэн\\
Үнэлгээ, андуурлын шинжилгээ, хугацаа, бичлэгийн туршилт, Gradio интерфэйс & гүйцэтгэсэн\\
Roboflow 984 зураг нэмж илрүүлэгчийг дахин сургах & үлдсэн (API түлхүүр)\\
Зогсоолын камерын бодит бичлэг цуглуулж шошголох, домэйн тохируулга & үлдсэн\\
Синтетикт ховор үсгийг тэнцвэржүүлэх, test-time augmentation & үлдсэн\\
Сурсан мөр хуваалт эсвэл хоёр ангит илрүүлэгч & үлдсэн\\
CTC beam search + дүрмийн хязгаарлалттай декод & үлдсэн\\
ONNX/TensorRT, INT8 экспорт, edge төхөөрөмж дээр хэмжих & үлдсэн\\
Орох-гарах бүртгэл, засварлах зай $\le1$ тулгалт, хаалга нээх дүрэм, операторын засварын лог & үлдсэн\\\bottomrule
\end{tabular}\end{table}
Зогсоолын системд нэгтгэх санал: илрүүлэгчийг кадр бүрт, уншигчийг зөвхөн мөр үүссэн үед ажиллуулах; $\ge3$ кадр, дүрэмд нийцсэн, margin $\ge0.8$ үед хаалга нээх; хамгийн сайн тайралт, санал нэгдсэн дугаар, margin-ыг бүртгэлийн санд хадгалж гарах үед нээлттэй сессүүдтэй засварлах зай $\le1$-ээр тулгах; татгалзсан, бага margin-тай уншилт бүрийг оператор засаж, засвар нь дараагийн сургалтын өгөгдөл болно. Хэдэн зуун засварласан дугаар хуримтлагдсаны дараа уншигчийг дахин тохируулж, тестийн олонлогийг 500-аас дээш дугаар болгосноор 95 хувийн мэдэгдэл $\pm2$ нэгж хувийн интервалтай болно.
